\documentclass[a4paper]{article}

%------------------------------------------------
% الحزم
%------------------------------------------------

\usepackage{enumitem}
\usepackage[margin=2.2cm]{geometry}
\usepackage{amsmath,amssymb}
\usepackage{graphicx}
\usepackage{array}
\usepackage{fancyhdr}
\usepackage[table]{xcolor}
\usepackage{tikz}
\usetikzlibrary{positioning,arrows.meta,shapes.geometric,calc}
\usepackage[skins]{tcolorbox}
\usepackage{polyglossia}

%------------------------------------------------
% اللغات والخطوط
%------------------------------------------------

\setdefaultlanguage[
    calendar=gregorian,
    numerals=maghrib
]{arabic}

\setotherlanguage{english}

% إذا لم تكن الخطوط مثبتة في الجهاز، تُؤخذ من المجلد fonts/
\IfFontExistsTF{Amiri}{%
    \newfontfamily\arabicfont[Script=Arabic, Scale=1.15]{Amiri}
}{%
    \newfontfamily\arabicfont[
        Script=Arabic, Scale=1.15,
        Path=fonts/, Extension=.ttf,
        UprightFont=*-Regular, BoldFont=*-Bold,
        ItalicFont=*-Italic, BoldItalicFont=*-BoldItalic
    ]{Amiri}
}

\IfFontExistsTF{AlBattar}{%
    \newfontfamily\arabicfontsf[Script=Arabic, Scale=1.2]{AlBattar}
    % خط العناوين (يعمل حتى بعد استعمال \textenglish)
    \newfontfamily\khatt[Script=Arabic, Scale=1.2]{AlBattar}
}{%
    \newfontfamily\arabicfontsf[
        Script=Arabic, Scale=1.2, Path=fonts/,
        BoldFont=ae_AlBattar.ttf
    ]{ae_AlBattar.ttf}
    % خط العناوين (يعمل حتى بعد استعمال \textenglish)
    \newfontfamily\khatt[
        Script=Arabic, Scale=1.2, Path=fonts/,
        BoldFont=ae_AlBattar.ttf
    ]{ae_AlBattar.ttf}
}

%------------------------------------------------
% الألوان: لون رئيسي هادئ + لون تنبيه فقط
%------------------------------------------------

\definecolor{asasi}{HTML}{3D6A8E}   % اللون الرئيسي: للرسوم والمخططات
\definecolor{tanbih}{HTML}{C0603A}  % لون التنبيه: للتحذيرات والحالات السلبية
\definecolor{ramadi}{HTML}{5F6B73}  % رمادي للنصوص الثانوية والأسهم
\definecolor{dakin}{HTML}{263238}   % لون العناوين

% أسماء قديمة تُحوَّل إلى اللونين فقط
\colorlet{azraq}{asasi}
\colorlet{fayruzi}{asasi}
\colorlet{banafsaji}{asasi}
\colorlet{akhdar}{asasi}
\colorlet{ahmar}{tanbih}
\colorlet{burtuqali}{asasi}

%------------------------------------------------
% الصناديق
%------------------------------------------------

\tcbset{
    asl/.style={
        enhanced,
        arc=2mm,
        boxrule=0.6pt,
        colback=white,
        colframe=black!18,
        colbacktitle=black!5,
        coltitle=dakin,
        titlerule=0pt,
        left=8pt, right=8pt, top=5pt, bottom=5pt,
        fonttitle=\khatt\bfseries,
        before skip=8pt, after skip=8pt
    }
}

% أهداف الدرس
\newtcolorbox{ahdaf}{asl, title={أهداف الدرس}}

% تعريف (خط رئيسي رفيع على الجانب)
\newtcolorbox{taarif}[1][تعريف]{asl, title={#1},
    borderline east={2pt}{0pt}{asasi}}

% مثال
\newtcolorbox{mithal}[1][مثال]{asl, title={#1}}

% تنبيه (اللون الوحيد البارز في النص)
\newtcolorbox{intabih}[1][انتبه!]{asl, colback=tanbih!5, colframe=tanbih,
    colbacktitle=tanbih, coltitle=white, title={#1}}

% ركن أستاذ اللغة العربية
\newtcolorbox{arabiyya}[1][في قسم اللغة العربية]{asl, title={#1}}

% نص أمر يكتب للأداة
\newtcolorbox{amr}[1][ramadi]{enhanced, arc=1mm, boxrule=0pt,
    colback=black!4, borderline east={2.5pt}{0pt}{#1},
    left=8pt, right=10pt, top=4pt, bottom=4pt,
    before skip=6pt, after skip=6pt}

% الخلاصة
\newtcolorbox{khulasa}{asl, title={خلاصة الدرس}}

% عنوان القسم
\newcounter{qism}
\newcommand{\qism}[1]{%
    \stepcounter{qism}%
    \par\bigskip
    {\large\khatt\bfseries{\color{asasi}\arabic{qism}.}\ {\color{dakin}#1}}\par
    \vspace{-2pt}{\color{black!20}\rule{\linewidth}{0.6pt}}\par\nopagebreak[4]
    \smallskip}

% عنوان فرعي
\newcommand{\far}[1]{\par\smallskip
    {\color{dakin}\khatt\bfseries {\color{asasi}$\blacktriangleleft$}\ #1}\par\nopagebreak[4]}

% كلمة مهمة (خط عريض فقط)
\newcommand{\kalima}[2][]{\textbf{#2}}

%------------------------------------------------
% أنماط الرسوم
%------------------------------------------------

\newcommand{\nodecolor}{}

\tikzset{
    % كتابة النص العربي من اليمين إلى اليسار داخل العقد
    % (لون النص يُعطى بالمفتاح lawn=... وليس text=...)
    every node/.append style={
        execute at begin node={\nodecolor\everypar{\beginR}\beginR},
        execute at end node=\endR},
    lawn/.code={\def\nodecolor{\color{#1}}},
    khatwa/.style={draw=#1, fill=#1!12, very thick, rounded corners=6pt,
        minimum height=1.2cm, text width=2.3cm, align=center,
        font=\khatt\bfseries},
    khatwa/.default=azraq,
    sahm/.style={-{Stealth[length=3mm]}, very thick, color=ramadi},
    raqm/.style={circle, fill=#1, lawn=white, inner sep=2pt,
        font=\small\bfseries},
}

%------------------------------------------------
% إعداد الصفحة
%------------------------------------------------

\renewcommand{\footrulewidth}{1pt}
\renewcommand{\headrulewidth}{1pt}

\rhead{\khatt المدرسة العليا للأساتذة ببوزريعة}
\lhead{\khatt ملحقة بومرداس}

\cfoot{\thepage}
\lfoot{\khatt الأستاذ: د. بن عبد الرزاق يوسف}

\parindent 0pt
\setlength{\headsep}{15pt}
\setlength{\parskip}{4pt}

\pagestyle{fancy}

\setlist[enumerate,1]{label=\arabic*), itemsep=2pt, topsep=3pt}
\setlist[itemize]{itemsep=2pt, topsep=3pt}

% الجداول
\arrayrulecolor{black!22}
\renewcommand{\arraystretch}{1.35}

% رأس جدول ملون
\newcommand{\ras}[1]{\cellcolor{black!7}\color{dakin}\bfseries #1}

% عنوان الدرس
\newcommand{\unwan}[3]{%
    \begin{center}
        {\color{asasi}\rule{\linewidth}{1.2pt}}\par\medskip
        {\khatt\color{ramadi} #1}\par\smallskip
        {\Large\khatt\bfseries\color{dakin} #2}\par\smallskip
        {\khatt\color{ramadi} #3}\par\medskip
        {\color{asasi}\rule{\linewidth}{1.2pt}}
    \end{center}}

%================================================
% بداية الوثيقة
%================================================

\begin{document}

%------------------------------------------------
% العنوان
%------------------------------------------------

\unwan{درس الفصل الأول}
{اكتشاف الذكاء الاصطناعي التوليدي}
{مدخل بسيط لطلبة قسم اللغة العربية}

\begin{ahdaf}
في نهاية هذا الدرس، تستطيع أن:
\begin{itemize}
    \item تعرّف الذكاء الاصطناعي التوليدي بكلمات بسيطة.
    \item تذكر أنواع المحتوى الذي ينتجه.
    \item تفهم لماذا تتغير الإجابة عندما نغيّر صياغة الطلب.
    \item تتحقق من المعلومة قبل استعمالها في درسك.
\end{itemize}
\end{ahdaf}

%================================================
\qism{ما الذكاء الاصطناعي؟}
%================================================

نستعمل كل يوم أجهزة تبدو «ذكية»: الهاتف يقترح الكلمة التالية،
والخريطة تقترح أقصر طريق.
وراء هذه الخدمات برامج نسميها \kalima{الذكاء الاصطناعي}.

\begin{taarif}
\textbf{الذكاء الاصطناعي} فرع من علوم الحاسوب.
يجعل الحاسوب يقوم بمهام تحتاج عادة إلى ذكاء الإنسان،
مثل الفهم والتعلم والإجابة عن الأسئلة.
\end{taarif}

%================================================
\qism{ما الذكاء الاصطناعي التوليدي؟}
%================================================

كلمة «التوليدي» مأخوذة من الفعل «ولَّد»، أي أنشأ شيئًا جديدًا.
فالذكاء الاصطناعي التوليدي لا يبحث عن إجابة جاهزة فقط،
بل \kalima[akhdar]{ينتج محتوى جديدًا} لم يكن موجودًا بهذا الشكل.

\begin{taarif}
\textbf{الذكاء الاصطناعي التوليدي} نوع من الذكاء الاصطناعي
ينتج نصوصًا أو صورًا أو أصواتًا جديدة، انطلاقًا من طلب يكتبه المستخدم.
\end{taarif}

\begin{center}
\begin{tikzpicture}
    \node[khatwa=azraq] (a) {المستخدم\\يكتب طلبًا};
    \node[khatwa=banafsaji, left=1cm of a] (b) {الأداة\\تحلل الطلب};
    \node[khatwa=akhdar, left=1cm of b] (c) {محتوى\\جديد};
    \draw[sahm] (a) -- (b);
    \draw[sahm] (b) -- (c);
\end{tikzpicture}
\end{center}

\begin{mithal}[تشبيه يقرّب الفكرة]
تخيّل تلميذًا قرأ آلاف النصوص العربية.
بعد ذلك، صار قادرًا على كتابة نص جديد بأسلوبه.
الأداة التوليدية تشبه هذا التلميذ، لكنها قرأت كمية ضخمة جدًا من النصوص.
\end{mithal}

%================================================
\qism{ماذا ينتج الذكاء الاصطناعي التوليدي؟}
%================================================

ينتج الذكاء الاصطناعي التوليدي أنواعًا مختلفة من المحتوى:

\begin{center}
\begin{tikzpicture}[
    far/.style={draw=#1, fill=#1!12, very thick, rounded corners=5pt,
        text width=2.4cm, align=center, minimum height=1cm,
        font=\khatt\bfseries}]
    \node[circle, draw=azraq, fill=azraq!15, ultra thick,
        text width=2.6cm, align=center, font=\khatt\bfseries]
        (m) {الذكاء الاصطناعي التوليدي};
    \node[far=akhdar]    (n1) at (0,2.6)    {نصوص};
    \node[far=burtuqali] (n2) at (-4.2,1.2) {صور};
    \node[far=asasi]     (n3) at (-4.2,-1.2){أصوات};
    \node[far=banafsaji] (n4) at (0,-2.6)   {فيديو};
    \node[far=fayruzi]   (n5) at (4.2,-1.2) {جداول\\ومخططات};
    \node[far=ramadi]    (n6) at (4.2,1.2)  {برامج\\حاسوبية};
    \foreach \i in {1,...,6} \draw[very thick, ramadi] (m) -- (n\i);
\end{tikzpicture}
\end{center}

\begin{arabiyya}
\begin{itemize}
    \item \textbf{نص:} قصة قصيرة لتلاميذ السنة الأولى متوسط.
    \item \textbf{صورة:} رسم يوضح مشهدًا من نص القراءة.
    \item \textbf{صوت:} قراءة مسجلة لقصيدة مع ضبط الحركات.
    \item \textbf{جدول:} جدول يلخص أنواع الخبر مع أمثلة.
\end{itemize}
\end{arabiyya}

%================================================
\qism{أشهر الأدوات}
%================================================

توجد أدوات كثيرة، وكلها تقبل الكتابة باللغة العربية.
من أشهرها:

\begin{center}
\rowcolors{2}{azraq!6}{white}
\begin{tabular}{|>{\centering\arraybackslash}m{4cm}|
                >{\centering\arraybackslash}m{4cm}|}
\hline
\ras{الأداة} & \ras{الشركة المطوِّرة} \\
\hline
\textenglish{ChatGPT}    & \textenglish{OpenAI}    \\ \hline
\textenglish{Gemini}     & \textenglish{Google}    \\ \hline
\textenglish{Claude}     & \textenglish{Anthropic} \\ \hline
\textenglish{DeepSeek}   & \textenglish{DeepSeek}  \\ \hline
\textenglish{Qwen}       & \textenglish{Alibaba}   \\ \hline
\textenglish{Perplexity} & \textenglish{Perplexity AI} \\ \hline
\end{tabular}
\end{center}

إذا أعطينا الطلب نفسه لهذه الأدوات، فلن نحصل على الإجابة نفسها.
تختلف الإجابات في الوضوح والطول والتنظيم وعدد الأمثلة.
لذلك من المفيد أن \kalima{نقارن} بين أكثر من أداة.

%================================================
\qism{الطلب يغيّر الإجابة}
%================================================

نسمي النص الذي نكتبه للأداة \kalima{الطلب} أو \kalima{الأمر}،
وبالإنجليزية \textenglish{Prompt}.
كلما كان الطلب واضحًا ودقيقًا، كانت الإجابة أفضل.

\far{طلب عام}
\begin{amr}[ahmar]
اشرح النعت.
\end{amr}

\far{طلب دقيق}
\begin{amr}[akhdar]
اشرح النعت لتلميذ في السنة الأولى متوسط،
بلغة بسيطة، مع ثلاثة أمثلة من الحياة اليومية.
\end{amr}

الطلب الثاني يحدد \kalima[akhdar]{المستوى} و\kalima[akhdar]{اللغة}
و\kalima[akhdar]{عدد الأمثلة}.
لذلك تأتي إجابته أقرب إلى حاجة الأستاذ.

%================================================
\qism{التحقق من إجابات الذكاء الاصطناعي}
%================================================

الأداة قد تخطئ، وقد تكتب المعلومة الخاطئة بثقة كبيرة.
نسمي هذه الظاهرة \kalima[ahmar]{الهلوسة}.

\begin{intabih}
قد تنسب الأداة بيتًا شعريًا إلى غير قائله،
أو تعطي إعرابًا خاطئًا لكلمة.
الإجابة الجميلة ليست بالضرورة إجابة صحيحة.
\end{intabih}

نتحقق من المعلومة المهمة باتباع هذه الخطوات:

\begin{center}
\begin{tikzpicture}[node distance=6mm]
    \node[khatwa=azraq, text width=2cm]     (a) {اقرأ\\الإجابة};
    \node[khatwa=fayruzi, text width=2cm, left=of a]  (b) {حدد\\المعلومة};
    \node[khatwa=burtuqali, text width=2cm, left=of b](c) {ابحث في\\مصدر موثوق};
    \node[khatwa=banafsaji, text width=2cm, left=of c](d) {قارن};
    \node[khatwa=akhdar, text width=2cm, left=of d]   (e) {قرّر};
    \draw[sahm] (a)--(b); \draw[sahm] (b)--(c);
    \draw[sahm] (c)--(d); \draw[sahm] (d)--(e);
\end{tikzpicture}
\end{center}

\textbf{أمثلة لمصادر موثوقة:} المعاجم المعروفة، كتب النحو المعتمدة،
الدواوين المحققة، والكتاب المدرسي الرسمي.

%================================================
\qism{كيف يستفيد أستاذ اللغة العربية؟}
%================================================

\begin{arabiyya}
يمكن للأستاذ أن يستعمل الذكاء الاصطناعي التوليدي في:
\begin{itemize}
    \item اقتراح أمثلة جديدة لقاعدة نحوية.
    \item إعداد أسئلة لفهم نص القراءة.
    \item تبسيط شرح درس في البلاغة.
    \item اقتراح موضوعات للتعبير الكتابي.
\end{itemize}
لكن الأستاذ يبقى هو \textbf{المسؤول}:
يقرأ، ويصحح، ويختار ما يناسب تلاميذه.
\end{arabiyya}

%================================================
\qism{اختبر فهمك}
%================================================

اكتب \kalima[akhdar]{صحيح} أو \kalima[ahmar]{خطأ}:

\begin{enumerate}
    \item كل ما يقدمه الذكاء الاصطناعي صحيح دائمًا.
    \item يمكن للذكاء الاصطناعي التوليدي إنتاج الصور.
    \item يمكن أن تتغير الإجابة حسب صياغة الطلب.
    \item يجب التحقق من المعلومات المهمة قبل استعمالها.
\end{enumerate}

{\small\color{ramadi} الإجابات: 1) خطأ \quad 2) صحيح \quad 3) صحيح \quad 4) صحيح}

\begin{khulasa}
\begin{itemize}
    \item الذكاء الاصطناعي التوليدي \textbf{ينتج محتوى جديدًا}:
          نصوصًا وصورًا وأصواتًا.
    \item جودة الإجابة مرتبطة بجودة \textbf{الطلب}.
    \item الأداة قد تخطئ، لذلك \textbf{نتحقق} دائمًا من مصدر موثوق.
    \item الذكاء الاصطناعي \textbf{مساعد} للأستاذ، وليس بديلًا عنه.
\end{itemize}
\end{khulasa}

\end{document}
